% Complete proofs for the selected 9 October 2001 List--Pettit preprint.

\begin{annexitem}{annex:p17.ja-models}{thm:p17-ja-models}{Valuations, weak consistency and agenda-relative closure}
\begin{proof}
For a valuation \(v\), put \(J_v=\{\varphi\in X:v\models\varphi\}\).
Opposing formulas have opposite truth values. The agenda condition therefore
makes \(J_v\) complete and weakly consistent. If \(J_v\models\varphi\),
then \(v\), which satisfies every member of \(J_v\), satisfies \(\varphi\).
For \(\varphi\in X\), this places \(\varphi\) in \(J_v\), proving closure.

Conversely, first suppose only that \(J\) is weakly consistent and closed.
If \(X\) is empty, then \(J\) is empty and is satisfiable. Otherwise the
opposing-proposition condition supplies a literal pair
\(\alpha,\neg\alpha\in X\). If \(J\) had no satisfying valuation, semantic
consequence would give both \(J\models\alpha\) and \(J\models\neg\alpha\):
with no model of \(J\), neither consequence has a countermodel. Closure would
put both in \(J\), contradicting weak consistency. Thus \(J\) has a model \(v\).

Now also assume completeness. We already have \(J\subseteq J_v\). Suppose
\(\varphi\in J_v\setminus J\). Completeness supplies either
\(\neg\varphi\in J\), or \(\psi\in J\) with \(\varphi=\neg\psi\).
In either case \(v\) must satisfy both opposing propositions, a contradiction.
Consequently \(J=J_v\).

Finally, \(\{p,q,\neg(p\land q)\}\) contains no formula together with its
literal negation. A valuation satisfying \(p\) and \(q\) satisfies their
conjunction and cannot satisfy its negation, so this set has no model. This
also shows why weak consistency without closure is insufficient.
\end{proof}
\end{annexitem}

\begin{annexitem}{annex:p17.ja-paradox}{thm:p17-ja-paradox}{The two- and three-premise profiles}
\begin{proof}
Each individual judgment set in the displayed three-person profile is the
restriction of a valuation to \(X_0\), so
Theorem~\ref{thm:p17-ja-models} gives all three rationality properties.
Individuals 1 and 2 accept \(p\), individuals 1 and 3 accept \(q\), and
individuals 2 and 3 accept \(\neg(p\land q)\). Only individual 1 accepts
\(p\land q\); only individual 3 accepts \(\neg p\); only individual 2
accepts \(\neg q\). Strict majority therefore returns precisely
\(\{p,q,\neg(p\land q)\}\).

Exactly one proposition in each of the three opposing pairs is selected, so the
set is complete and weakly consistent. It has no model by
Theorem~\ref{thm:p17-ja-models}, and it omits the consequence \(p\land q\).
Because an unsatisfiable set semantically entails every formula, its
agenda-relative closure is \(X_0\). In particular, the closure contains both
\(p\land q\) and its negation.

For the three-premise profile, use the agenda consisting of \(p,q,r\), their
negations, and \(p\land q\land r\) and its negation. Each row again extends to
a valuation. The respective premise-support coalitions are \(\{1,3\}\),
\(\{2,3\}\) and \(\{1,2\}\); each has two members. Their intersection is
empty. Each individual falsifies one conjunct and therefore falsifies the
conclusion, whose negation has three supporters. This proves the claimed
unanimous rejection despite separate premise majorities.
\end{proof}
\end{annexitem}

\begin{annexitem}{annex:p17.ja-counts}{thm:p17-ja-counts}{From one systematic Boolean function to a count function}
\begin{proof}
Condition S supplies one function \(f\) on Boolean \(n\)-vectors. To derive
its symmetry from A, it is necessary to know that the vectors being compared
actually occur in admissible profiles. Fix any \(d=(d_1,\ldots,d_n)\).
For each individual choose a valuation with \(v_i(p)=d_i\), assign any fixed
values, for example zero, to all other atoms, and put
\(J_i=\{\varphi\in X:v_i\models\varphi\}\).
These are rational judgments by Theorem~\ref{thm:p17-ja-models}, so U admits
the profile. The support vector of \(p\) is exactly \(d\).

For a permutation \(\sigma\), the permuted profile has support vector
\((d_{\sigma(1)},\ldots,d_{\sigma(n)})\) on \(p\).
Anonymity gives the same collective membership of \(p\) at these two
profiles. Applying S at both profiles proves
\[
 f(d_1,\ldots,d_n)=f(d_{\sigma(1)},\ldots,d_{\sigma(n)}).
\]
Any two Boolean vectors with equally many ones are related by a permutation:
match the positions containing ones bijectively, and do the same for the zero
positions. Therefore \(f\) is constant on each support-count class. Define
\(g(k)\) as its value on a vector with exactly \(k\) ones. All integers
\(0\leq k\leq n\) have such a vector, and the preceding argument makes
this definition independent of the chosen vector. Applying it to the support
vector of any proposition proves the assertion.
\end{proof}
\end{annexitem}

\begin{annexitem}{annex:p17.ja-impossibility}{thm:p17-ja-impossibility}{The original even- and odd-electorate impossibility construction}
\begin{proof}
Suppose \(F:\mathcal J_X^n\to\mathcal J_X\) satisfies U, A and S. Write
\(c=p\land q\), so \(p,q,c,\neg c\in X\), and use the count function
\(g\) from Lemma~\ref{thm:p17-ja-counts}.

Construct valuations on \(p,q\) by the following table, using only rows whose
index is at most \(n\):
\[
 \begin{array}{c|cccc}
 i& p&q&c&\neg c\\ \hline
 1&1&1&1&0\\
 2&1&0&0&1\\
 3&0&1&0&1\\
 i>3,\ i\text{ even}&1&1&1&0\\
 i>3,\ i\text{ odd}&0&0&0&1
 \end{array}
\]
Values in the last two columns follow from the first two; they are not
independent votes. Assign zero to every other atom in every row, and take
each \(J_i\) to be the corresponding truth set restricted to the entire
agenda \(X\). Theorem~\ref{thm:p17-ja-models} makes every input complete,
weakly consistent and closed, and U admits the profile. Let
\(J=F(J_1,\ldots,J_n)\).

\emph{Even case.} If \(n=2\), exactly one person accepts \(c\), and exactly
one accepts \(\neg c\). For even \(n\geq4\), the first four rows give two
supporters of each, and every subsequent pair of rows adds one supporter of
each. Thus in every even case
\[
 s_c=s_{\neg c}=n/2.
\]
The shared count function makes \(c\in J\) equivalent to \(\neg c\in J\).
Completeness of \(J\) requires at least one of this pair, so both must belong
to \(J\). This violates weak consistency. Equivalently, if the common value
were zero, completeness itself would fail.

\emph{Odd case.} Write \(n=2m+1\). Since \(n\geq2\), here \(m\geq1\).
The first three rows give two supporters each to \(p,q,\neg c\) and one to
\(c\). The remaining rows comprise \(m-1\) pairs, each containing a
\((1,1)\) row and a \((0,0)\) row. Each pair adds one supporter to each of
the four formulas. Consequently
\[
 s_p=s_q=s_{\neg c}=m+1,\qquad s_c=m.
\]
Systematicity and anonymity therefore give either all or none of
\(p,q,\neg c\) in \(J\).

If all three belong to \(J\), its two premises entail \(c\). Because
\(c\in X\), closure puts \(c\) in \(J\), contradicting the simultaneous
membership of \(\neg c\).

If none belongs to \(J\), completeness applied to the pair \(c,\neg c\)
puts \(c\) in \(J\). Since \(c\models p\) and \(p\in X\), closure then
puts \(p\) in \(J\), contradicting its assumed absence. These two cases
exhaust \(g(m+1)\in\{0,1\}\).

Every finite \(n\geq2\) is even or odd, and both cases contradict the proposed
output requirements. This is the two-step construction of the selected source's
appendix, with the realization of all input profiles and the case \(n=2\)
made explicit \cite[pp.~27--29]{ListPettit2001JudgmentPreprint}.
For a proposed rule on a larger domain containing every rational profile,
restrict it to \(\mathcal J_X^n\). Anonymity, the same systematic Boolean
function and the required rational outputs survive restriction. This would
give the impossible rule just excluded, proving the larger-domain assertion.
\end{proof}
\end{annexitem}

\begin{annexitem}{annex:p17.ja-median}{thm:p17-ja-median}{Majority equals the median on an aligned profile}
\begin{proof}
Take an order \(i_1,\ldots,i_{2m+1}\) witnessing alignment and let
\(i_{m+1}\) be its middle member. Fix a proposition \(\varphi\).
If its supporters form an initial segment containing the middle member,
that segment has at least \(m+1\) members. If the initial segment does not
contain the middle member, it has at most \(m\). For a final segment the
same counts hold in the reverse direction. Empty and full segments also
satisfy these statements. Therefore
\[
 s_\varphi>m
 \quad\Longleftrightarrow\quad
 \varphi\in J_{i_{m+1}}.
\]
Since this holds for every \(\varphi\in X\), the majority output is exactly
\(J_{i_{m+1}}\), which is rational by the input assumption.

The admissible domain is invariant under relabeling individuals: permuting
the labels in a witnessing order gives another witnessing order. Majority
counts themselves are invariant under such relabeling, proving A. The fixed
function \(f(d)=1\) when \(\sum_i d_i>m\) proves S on this domain. There is
no claim of U over all rational profiles.

For the displayed five-person example, the supporters of \(p\) are the
initial three positions, and those of \(q\) the final two. No one accepts
\(p\land q\), and everyone accepts its negation. The opposing atoms have
the complementary segments. Majority thus selects \(p,\neg q,\neg(p\land q)\),
which are exactly the judgments of the middle individual, person 2. The result
does not require naming one person as median in every profile.
\end{proof}
\end{annexitem}

\begin{annexitem}{annex:p17.ja-authority-priority}{thm:p17-ja-authority-priority}{Authority, independent premises and missing completeness}
\begin{proof}
For the fixed-authority rule \(F_j(\mathbf J)=J_j\), the output is a rational
input set. The rule is defined on every rational profile and is systematic
with \(f(d)=d_j\). To refute anonymity, choose another individual
\(h\ne j\), let \(J_j\) arise from a valuation making \(p\) true and
\(J_h\) from one making \(p\) false. Swapping these individuals changes
the output's membership of \(p\). Both profiles are admissible.

For the priority rule, the votes and its fixed anonymous tie convention assign
exactly one value to each member of \(Y\). The stipulated independence of
\(Y\) provides a valuation \(v\) satisfying all the signed choices in
\(B_Y\). Consequently every formula in \(\operatorname{cl}_X(B_Y)\) is
true in \(v\). The collective set is satisfiable and cannot contain an
opposing pair.

It is closed because consequences of consequences are consequences of the
original set: if \(\operatorname{cl}_X(B_Y)\models\varphi\), then every
model of \(B_Y\) satisfies \(\operatorname{cl}_X(B_Y)\) and hence
\(\varphi\). When \(\varphi\in X\), this places it in the displayed
closure. The input choices exist for every rational profile, proving U, and
permuting individuals does not change the votes or anonymous tie decisions,
proving A.

Suppose the assignment to \(Y\) determines every agenda formula. For every
\(\varphi\in X\), all valuations satisfying the chosen signs agree on its
truth value. If they make it true, it belongs to the closure. If they make it
false, the appropriate opposite formula supplied by the agenda condition is
true in all those valuations and belongs to the closure. This proves
completeness. Equivalently, the closure is the truth set on \(X\) of any one
of its satisfying valuations.

For failure of completeness without this condition, take \(Y=\{p\land q\}\)
on \(X_0\). This singleton is logically independent: its true and false
assignments both have models. At the doctrinal profile the conclusion vote
selects \(\neg(p\land q)\). Its three \((p,q)\)-models are
\((0,0),(1,0),(0,1)\). They agree only on the negated conjunction among
the six agenda formulas, so the closure is precisely
\(\{\neg(p\land q)\}\), missing an answer on both atomic issues.

For failure of S even when the prioritized questions determine the agenda,
take \(Y=\{p,q\}\) and compare these two three-person valuation profiles:
\[
 \begin{array}{c|cc|cc}
 &\multicolumn{2}{c|}{\text{profile A}}&\multicolumn{2}{c}{\text{profile B}}\\
 i&p&q&p&q\\ \hline
 1&1&1&1&1\\
 2&1&0&0&0\\
 3&0&1&0&0
 \end{array}
\]
The individual support vector for \(c=p\land q\) is \((1,0,0)\) in both.
At A both premises have majority support, so priority followed by closure
accepts \(c\). At B both premises have majority rejection, so it rejects
\(c\). No function of the votes on \(c\) alone can give these two different
answers. This violates the independence component of S, not just the demand
to use identical rules for different propositions.
\end{proof}
\end{annexitem}

\begin{annexitem}{annex:p17.ja-unanimity}{thm:p17-ja-unanimity}{Intersection of rational judgment sets}
\begin{proof}
Let \(J=\bigcap_iJ_i\). Since \(n\geq2\), it is a subset of at least one
individual rational judgment set and has any model of that set as a model.
It is therefore satisfiable and weakly consistent.

If \(J\models\varphi\) with \(\varphi\in X\), each \(J_i\), being a
superset of \(J\), also entails \(\varphi\). Each is closed, so
\(\varphi\in J_i\) for all \(i\), whence \(\varphi\in J\).
Thus \(J\) is already closed. Intersection is defined on every rational
profile and is invariant under permutations; these give U and A. Its membership
condition is exactly that every input bit be 1, giving S.

Choose rational inputs disagreeing on \(p\). Then neither \(p\) nor
\(\neg p\) belongs to their intersection. Thus completeness is not guaranteed.
\end{proof}
\end{annexitem}

\begin{annexitem}{annex:p17.ja-supermajority}{thm:p17-ja-supermajority}{Common support, strict thresholds and closure}
\begin{proof}
For rational inputs, the support sets of the two formulas in any issue pair
are complementary. Since \((k-1)/k\geq1/2\) for \(k\geq2\), two opposite
formulas cannot both receive strictly more than \((k-1)n/k\) votes. There
are therefore at most \(k\) formulas in \(H_k(\mathbf J)\).

List its accepted formulas as \(\varphi_1,\ldots,\varphi_\ell\), where
\(0\leq\ell\leq k\). For each accepted formula,
\[
 \#(N\setminus N_{\varphi_j})<n/k.
\]
If \(\ell\geq1\), the elementary finite union bound gives
\[
 \#\bigcup_{j=1}^{\ell}(N\setminus N_{\varphi_j})
 \ \leq\ \sum_{j=1}^{\ell}\#(N\setminus N_{\varphi_j})
 \ <\ \ell n/k\ \leq\ n.
\]
Some individual lies outside this union, hence accepts every selected formula.
The rational judgment set of that individual has a model by
Theorem~\ref{thm:p17-ja-models}; that model satisfies the whole collective
set. If \(\ell=0\), the empty set is satisfiable directly. This proves the
claimed consistency guarantee without assuming that any one proposition
with subthreshold support must be false.

The support test is defined for every rational profile and depends only on the
number of ones in the proposition's vote vector. It consequently satisfies U,
A and S. Taking its agenda-relative consequences preserves a satisfying model,
and the consequence-of-consequences argument from
Annex~\ref{annex:p17.ja-authority-priority} proves closure. The resulting rule
still has U and A, since neither its domain nor its symmetry has changed.

The remaining limitations can all be witnessed on \(X_0\), with \(k=3\) and
four individuals. The threshold is strictly above \(8/3\), hence requires
three votes. Compare
\[
 \begin{array}{c|cc|cc}
 &\multicolumn{2}{c|}{\text{profile A}}&\multicolumn{2}{c}{\text{profile B}}\\
 i&p&q&p&q\\ \hline
 1&1&1&1&1\\
 2&1&1&1&1\\
 3&1&0&0&0\\
 4&0&1&0&0
 \end{array}
\]
At A, \(p\) and \(q\) have three supporters each, while their conjunction
and its negation each have two. The negated atoms have one supporter each.
Thus \(H_3\) returns \(\{p,q\}\), which is incomplete and not closed.
Its closure is \(\{p,q,p\land q\}\). At B, every agenda formula has two
supporters, so \(H_3\) is empty. No formula of \(X_0\) is a tautology;
its closure is therefore also empty and incomplete.

The support vector of \(p\land q\) is \((1,1,0,0)\) in both profiles, but
\(G_3\) accepts it at A and not at B. This disproves systematicity for the
closure rule. Finally, replacing the strict threshold by at least \(2/3\)
for three individuals and three issues admits \(p,q,\neg(p\land q)\) at
the doctrinal profile. Their joint unsatisfiability and inconsistent closure
were proved in Annex~\ref{annex:p17.ja-paradox}. Thus the strict boundary
cannot be weakened in the stated general guarantee.
\end{proof}
\end{annexitem}

\begin{annexitem}{annex:p17.ja-boundaries}{thm:p17-ja-boundaries}{What changes when a requirement is removed}
\begin{proof}
On rational inputs, the two sides of every opposing pair receive complementary
votes. With odd \(n\), precisely one side has a strict majority. The raw
majority set is therefore complete and weakly consistent. Its count rule
satisfies A and S and is defined on every rational profile, giving U. The
doctrinal example shows that closure and satisfiability can nevertheless fail.

Closing that example yields \(X_0\), with opposing pairs. Closure can also
destroy S: use the two profiles A and B in
Annex~\ref{annex:p17.ja-authority-priority}. At A the majority set is the
unsatisfiable doctrinal set, so its closure contains \(p\land q\). At B the
majority set is \(\{\neg p,\neg q,\neg(p\land q)\}\), already a rational
set, and its closure does not contain \(p\land q\). The individual vote
vector for that conjunction is \((1,0,0)\) in both cases. Thus the closure
operation is not guaranteed to preserve systematicity.

The constant-full rule returns every agenda proposition, so it is complete
and closed by definition. It is defined on all profiles, permutation-invariant
and systematic with \(f(d)=1\) for every vector. It contains both \(p\) and
\(\neg p\), so it is weakly inconsistent. This example demonstrates the
logical effect of dropping consistency; it supplies no responsiveness or
normative justification beyond the formal conditions actually stated.

If the electorate has one individual, the identity rule returns that rational
input. U is immediate, A has only the identity permutation to check, and
\(f(d)=d\) gives S. This is outside the hypothesis \(n\geq2\).

Finally, consider an agenda consisting only of independent atoms and their
negations, with an odd electorate. Assign to each of its atoms the majority
truth value. Extend this assignment to the remaining propositional atoms, for
example by assigning them zero. The resulting valuation has exactly the
majority judgment set on this agenda. The same direct truth-set argument as
in Annex~\ref{annex:p17.ja-models} proves completeness, consistency and closure.
The dependent conjunction required for the impossibility theorem is absent.
\end{proof}
\end{annexitem}
